\documentclass[12pt]{article}
\usepackage[utf8]{inputenc}
\usepackage[T1]{fontenc}
\usepackage{hyperref}
\title{Gbin Language Resources for K-12 Curriculum Development}
\author{Compiled by Claude Sonnet 5.5 (claude-sonnet-5-5) for Andr\'as Kornai}
\date{}
\begin{document}
\maketitle

\section{Introduction}
Gbin (Gb\~i; ISO 639-3 xgb) is an extinct language of C\^ote d'Ivoire, reported around Bondoukou. It is usually described as Mande. Sources disagree on its placement. Paperno groups it with Beng in Southern Mande, while Glottolog lists it as unclassified within the Nwa-Ben group of Southeastern Mande. Searches found almost no teaching or revitalization material. Nearly everything known about the language comes from one 1904 source, Delafosse (1904).

\section{Primary Educational Resources}
No primer, textbook, school programme or curriculum for Gbin was found. No native-language teaching materials exist in the search results.

\section{Linguistic and Reference Materials}
\begin{itemize}
\item Wikipedia, ``Gbin language'', a short overview: \url{https://en.wikipedia.org/wiki/Gbin_language}
\item Glottolog entries for neighbouring Mande languages, useful for classification only: \url{https://glottolog.org/resource/languoid/id/kono1267} and \url{https://glottolog.org/resource/languoid/id/samo1309}. These are not Gbin pages.
\item An article in \emph{Mandenkan} on the Beng and Gbin area: \url{https://isidore.science/document/10.4000/mandenkan.555}. I saw only a search snippet, so its exact relevance to Gbin is unverified.
\item Delafosse (1904) is the principal attestation of Gbin. No URL for it was seen in the search results.
\end{itemize}

\section{Religious and Community Resources}
No scripture translations or community organisations for Gbin were found.

\section{Digital and Audio Resources}
No audio, video, corpus or software resources were found.

\section{Curriculum Development Context}
There are no known speakers and the documentation is minimal, so a curriculum cannot use the language for much. A K-12 heritage unit could at most present the history of the Gbin people, their relation to the neighbouring Beng, and what is known of the language. Any word material would have to be taken from Delafosse (1904), after checking it directly. Work with Beng communities and Ivorian linguists would be needed.
\end{document}